\subsection*{2.2.}
Let now $f:X\longrightarrow\Spec(\mathbb F_q)$ be a proper scheme over
$\mathbb F_q$. Let $k$ be an algebraic closure of $\mathbb F_q$, and put

\begin{equation}\tag{2.2.1}
\bar X=X\times_{\mathbb F_q}k.
\end{equation}

The étale sheaf $R^if_*(\mathbb Z/p\mathbb Z)$ on
$\Spec(\mathbb F_q)$ identifies with
$H^i_{\mathrm{et}}(\bar X,\mathbb Z/p\mathbb Z)$, viewed as a
$\Gal(k/\mathbb F_q)$-module. Let

\begin{equation}\tag{2.2.2}
\mathfrak F_q:H^i(X,\mathcal O_X)\longrightarrow H^i(X,\mathcal O_X)
\end{equation}

denote the linear operation (it is $q$-linear, but the ground field is
$\mathbb F_q$) induced by

\begin{equation}\tag{2.2.3}
\mathfrak F_q:\mathcal O_X\longrightarrow\mathcal O_X,
\qquad \mathfrak F_q(f)=f^q.
\end{equation}

\subsubsection*{Proposition 2.2.4.}
Let $\varphi_q\in\Gal(k/\mathbb F_q)$ be the ``Frobenius'' element
$\varphi_q(\lambda)=\lambda^q$, and let $i\in\mathbb Z$. There is an
identity of characteristic polynomials

\begin{equation}\tag{2.2.4.1}
\det_{\mathbb F_p}\!\left(
1-t\varphi_q^{-1}\mid H^i_{\mathrm{et}}(\bar X,\mathbb Z/p\mathbb Z)
\right)
=
\det_{\mathbb F_q}\!\left(
1-t\mathfrak F_q\mid H^i(X,\mathcal O_X)
\right).
\end{equation}

Here the left-hand side is the characteristic polynomial of the element
$\varphi_q^{-1}$ of $\pi_1(\Spec(\mathbb F_q))$ acting on the Galois module
$H^i_{\mathrm{et}}(\bar X,\mathbb Z/p\mathbb Z)$.

Proposition~2.2.4 is a corollary of the following result.

\subsubsection*{Proposition 2.2.5.}
Let $X$ be proper over $\Spec(\mathbb F_q)$. There is a canonical
isomorphism

\begin{equation}\tag{2.2.6}
H^i(X,\mathcal O_X)_{\mathrm{ss}}\otimes_{\mathbb F_q}k
\simeq
H^i_{\mathrm{et}}(\bar X,\mathbb Z/p\mathbb Z)
\otimes_{\mathbb F_p}k
\end{equation}

which carries the automorphism induced by $\mathfrak F_q$ on the left-hand
side to the automorphism induced by $\varphi_q^{-1}$ on the right-hand side.
The latter may also be interpreted as
$(\mathrm{Fr}_q^X)^*\otimes_{\mathbb F_p}k$, where
$\mathrm{Fr}_q^X:\bar X\longrightarrow\bar X$ is the Frobenius
$k$-endomorphism (cf. SGA~5, XVIII~2.3.4).

\medskip
\noindent\emph{Proof.}
Since

\begin{equation}\tag{2.2.7}
H^i(\bar X,\mathcal O_{\bar X})\simeq
H^i(X,\mathcal O_X)\otimes_{\mathbb F_q}k,
\end{equation}

Proposition~1.0.10 gives

\begin{equation}\tag{2.2.8}
H^i(\bar X,\mathcal O_{\bar X})_{\mathrm{ss}}\simeq
H^i(X,\mathcal O_X)_{\mathrm{ss}}\otimes_{\mathbb F_q}k.
\end{equation}

Thus the isomorphism (2.2.6) is obtained by composing (2.2.8) with
(2.1.1).

In view of (2.2.7), the $q$-linear operation

\begin{equation}\tag{2.2.9}
\mathfrak F_q:H^i(X,\mathcal O_X)\longrightarrow H^i(X,\mathcal O_X)
\end{equation}

has a $q$-linear extension (cf. 1.0.5)

\begin{equation}\tag{2.2.10}
\overline{\mathfrak F}_q:H^i(\bar X,\mathcal O_{\bar X})
\longrightarrow H^i(\bar X,\mathcal O_{\bar X}),
\end{equation}

and also a $k$-linear extension, again denoted by
$\overline{\mathfrak F}_q$,

\begin{equation}\tag{2.2.11}
\overline{\mathfrak F}_q:H^i(\bar X,\mathcal O_{\bar X})
\longrightarrow H^i(\bar X,\mathcal O_{\bar X}).
\end{equation}

Let

\begin{equation}\tag{2.2.12}
\varphi_q:H^i(\bar X,\mathcal O_{\bar X})
\longrightarrow H^i(\bar X,\mathcal O_{\bar X})
\end{equation}

denote the $q$-linear operation whose fixed points are
$H^i(X,\mathcal O_X)$. Under (2.2.7), the operators
$\mathfrak F_q$, $\varphi_q$, and $\overline{\mathfrak F}_q$ may also be
written respectively as
$\mathfrak F_q\otimes\mathrm{id}_k$,
$\mathrm{id}\otimes\varphi_q$, and
$\mathfrak F_q\otimes\varphi_q$. They therefore commute with one another,
and

\begin{equation}\tag{2.2.13}
\overline{\mathfrak F}_q
=\varphi_q\mathfrak F_q
=\mathfrak F_q\circ\varphi_q.
\end{equation}

They all commute with $\mathfrak F_p$, so $\mathfrak F_q$ and $\varphi_q$
act on both
$H^i(\bar X,\mathcal O_{\bar X})^{1-\overline{\mathfrak F}_q}$ and
$H^i(\bar X,\mathcal O_{\bar X})^{1-\mathfrak F_p}$. By (2.2.13),
$\mathfrak F_q$ and $\varphi_q$ are inverse to one another on
$H^i(\bar X,\mathcal O_{\bar X})^{1-\overline{\mathfrak F}_q}$. Hence the
action of $\mathfrak F_q$ on

\[
H^i(\bar X,\mathcal O_{\bar X})_{\mathrm{ss}}
\simeq
H^i(\bar X,\mathcal O_{\bar X})^{1-\overline{\mathfrak F}_q}
\otimes_{\mathbb F_q}k
\qquad\text{(cf. 1.1.11)}
\]

is inverse to the action of $\varphi_q\otimes\mathrm{id}_k$ on

\[
H^i(\bar X,\mathcal O_{\bar X})^{1-\overline{\mathfrak F}_q}
\otimes_{\mathbb F_q}k.
\]

Writing

\begin{equation}\tag{2.2.14}
H^i(\bar X,\mathcal O_{\bar X})^{1-\overline{\mathfrak F}_q}
\simeq
H^i(\bar X,\mathcal O_{\bar X})^{1-\mathfrak F_p}
\otimes_{\mathbb F_p}\mathbb F_q,
\end{equation}

the automorphism $\varphi_q$ of the left-hand side of (2.2.14) becomes,
via (2.2.14), the automorphism induced by $\varphi_q$ on
$H^i(\bar X,\mathcal O_{\bar X})^{1-\mathfrak F_p}
\otimes_{\mathbb F_p}\mathbb F_q$. Consequently the automorphism induced by
$\mathfrak F_q$ on
$H^i(\bar X,\mathcal O_{\bar X})_{\mathrm{ss}}$ becomes, via (2.1.1), the
automorphism induced by $\varphi_q^{-1}$ on
$H^i_{\mathrm{et}}(\bar X,\mathbb Z/p\mathbb Z)
\otimes_{\mathbb F_p}k$. This proves Proposition~2.2.5.

\medskip
\noindent\emph{Source note.}
The final printed sentence of the proof says $\varphi_q$ where the proposition
and the immediately preceding inverse relation require $\varphi_q^{-1}$.
The English follows the mathematically coherent reading and the corrected
French workpass; the diplomatic French retains the printed wording.
